
\textbf{Proposition 4.5. (Comportement fonctoriel).} \emph{Soit @@M0@@ un morphisme entre deux @@M1@@-sch\'emas propres, et soit @@M2@@ un module inversible sur @@M3@@. Alors}~:
\begin{enumerate}[label=(\roman*)]
\item \emph{Si @@M4@@ est @@M5@@-\'equivalent \`a @@M6@@ (resp. num\'eriquement \'equivalent \`a @@M7@@), alors @@M8@@ l'est.}
\item \emph{R\'eciproquement, lorsque @@M9@@ est surjectif, si @@M10@@ est @@M11@@-\'equivalent \`a @@M12@@ (resp. num\'eriquement \'equivalent \`a @@M13@@), alors @@M14@@ l'est.}
\end{enumerate}

En effet, l'assertion (i) (resp. (ii)) pour la @@M15@@-\'equivalence r\'esulte aussit\^ot de 4.1 (resp. 4.3). D'autre part, soient @@M16@@ une courbe int\`egre ferm\'ee dans @@M17@@, et @@M18@@ son image, munie de la structure r\'eduite induite. Lorsque @@M19@@, on a la formule de projection (X 4.4.4)
@@M20@@
o\`u @@M21@@ (resp. @@M22@@) est le corps de fonctions de @@M23@@ (resp. @@M24@@), et lorsque @@M25@@, on a @@M26@@. Par cons\'equent, si @@M27@@ est num\'eriquement \'equivalent \`a @@M28@@, alors @@M29@@\,; d'o\`u (i).

R\'eciproquement, lorsque @@M30@@ est surjective, \'etant donn\'ee une courbe int\`egre ferm\'ee @@M31@@ dans @@M32@@, on peut trouver une courbe int\`egre ferm\'ee @@M33@@ dans @@M34@@ telle que @@M35@@~: soient @@M36@@ le point g\'en\'erique de @@M37@@, @@M38@@ un point de la fibre @@M39@@ alg\'ebrique sur @@M40@@, et @@M41@@ l'adh\'erence de @@M42@@. Alors 4.5.1 montre que @@M43@@ implique @@M44@@\,; d'o\`u (ii).

\textbf{Th\'eor\`eme 4.6. (La caract\'erisation num\'erique de @@M45@@).} \emph{Soient @@M46@@ un corps alg\'ebriquement clos, @@M47@@ un @@M48@@-sch\'ema propre, et @@M49@@ un faisceau inversible sur @@M50@@. Alors les conditions suivantes sont \'equivalentes}~:
\begin{enumerate}[label=(\roman*)]
\item \emph{@@M51@@ est @@M52@@-\'equivalent \`a @@M53@@.}
\item \emph{Pour tout faisceau coh\'erent @@M54@@ sur @@M55@@, on a}
